% =====================================================================
% PISTES POUR DE FUTURS MEMOIRES D'ACTUARIAT
%
% ANNEXE EXIGEE PAR L'INSTITUT DES ACTUAIRES. Le paragraphe 5.1.h des
% Recommandations Jury applicables depuis le 1er avril 2026 impose deux
% annexes, dont celle-ci : << Pistes eventuelles pour de futurs memoires
% actuariels >>, qui << peuvent reprendre de facon non limitative des pistes
% d'enrichissement du memoire presente mais aussi des axes d'etudes
% abandonnees et n'ayant pas conduit a une impasse definitive >>, et qui
% << constituent un legs a l'entreprise ayant accueilli l'etudiant mais aussi
% un apport pour la communaute actuarielle >>.
%
% CREE LE 21 SEPTEMBRE 2026. Aucune piste inventee : chacune est un travail
% deja instruit dans le memoire et laisse ouvert pour une raison ecrite. Les
% grandeurs chiffrees reprises ici viennent toutes d'ailleurs dans le memoire et
% citent le script qui les imprime : 66 % et 28,9 % du chapitre de la donnee
% manquante, et les elasticites 0,95 / 0,39 / 2,42 du script 74.
%
% DEUX PISTES DE STRUCTURE AJOUTEES LE 23 SEPTEMBRE 2026, sur demande de Kelian.
% Motif : sur les six pistes d'origine, CINQ demandaient de la donnee et une
% seule etait une evolution de modelisation. L'annexe se lisait donc comme
% << il manque des donnees a ce travail >>, ce qui est defensif. Les deux
% ajoutees portent sur la STRUCTURE du modele, aux deux endroits ou il se sait
% le plus contraint : l'additivite des couts et le perimetre d'une seule entite.
% Elles sont placees en fin, l'annexe etant ordonnee du plus executable au plus
% exploratoire, et l'encadre de cloture est mis a jour avec elles.
% =====================================================================

\chapter{Pistes pour de futurs mémoires d'actuariat}
\label{ann:pistes}

Les pistes suivantes prolongent des travaux préparés dans ce mémoire et restés ouverts.
Elles vont de la mise en œuvre d'un protocole existant à la construction d'un nouveau modèle.

\section{Exécuter l'élicitation structurée}

L'annexe~\ref{ch:elicitation} décrit un protocole de Cooke complet, préparé et diffusé, resté sans
réponse. Le questionnaire est rédigé ; son étalonnage comporte dix questions de graine, dont les
valeurs vraies sont calculées et contrôlées. Il manque le panel. Un futur mémoire partirait d'un
instrument prêt : il constituerait le panel, mesurerait la calibration des experts, pondérerait leurs
avis, puis confronterait les poids obtenus aux directions que le corpus indique déjà.

Les réponses aux questions de graine restent confidentielles afin de préserver l'étalonnage du questionnaire.

\section{Documenter la dépendance entre gouvernance et gestion des incidents}

La section~\ref{ch:livrable} hiérarchise les données manquantes par leur valeur d'information.
Fixer le seul sens de la dépendance entre la gouvernance du risque TIC et la gestion des incidents
lèverait $13{,}3\,\%$ de l'ambiguïté de capital en moyenne sur quatre graines, de $10{,}6$ à
$18{,}5\,\%$. Fixer celui du lien entre tests et prestataires tiers n'en lèverait que $2{,}3\,\%$.
Le corpus public ne couvre pas les deux dépendances les plus utiles. Un mémoire conduit chez un
assureur disposant d'une chronologie interne d'incidents, horodatée et classée par domaine de
contrôle, mesurerait de combien la bande se resserre réellement.

% SOURCES-SCRIPTS: 55 117

\section{Le second codage en aveugle du corpus}

La direction des dépendances est codée par un seul lecteur, limite déclarée. La réception d'un
second codage est outillée : le script d'accord ne produit rien sans le second jeu. Il dirait si la
direction retenue tient à la lecture ou au texte.

\section{Calibrer l'amplitude des liens, et non leur seul sens}

Même avec les dix directions connues, une bande de capital subsisterait, imputable à l'amplitude des
coefficients. Connaître le sens d'un lien ne donne pas sa force. Toute la valeur d'information
mesurée ici, $32{,}7\,\%$ de resserrement en moyenne si l'on retenait le sens documenté par les
rapports publics, porte sur la direction. Un mémoire qui estimerait les amplitudes, par exemple sur
des séries d'incidents d'une même entité où la sévérité transmise est observable, traiterait la moitié
du problème que celui-ci laisse intacte.

% SOURCES-SCRIPTS: 66 55 117

\section{Dérive d'échelle et passage à l'entité}

Le backtest couvre la forme de la queue et la loi de fréquence, mais pas le niveau de la charge
annuelle : une dérive est examinée, sans validation hors échantillon d'un modèle non stationnaire. Un mémoire sur
la non-stationnarité de la sévérité cyber partirait de ce diagnostic chiffré. De même, le besoin de
capital est ramené à l'entité par ses actifs et des lectures d'exposition, sans seuil de
validité statistique démontré. Un mémoire conduit dans une entité, avec son exposition et son registre de
prestataires, remplacerait cet ancrage par une mesure directe.

\section{Le coût d'un sinistre multi-piliers}

Lorsqu'un incident touche plusieurs piliers, le mémoire additionne leurs coûts. C'est la plus lourde
des hypothèses qu'il ne teste pas ; il la relâche par un exposant sur le nombre de piliers touchés
(\S\ref{sec:additivite-couts}). Sur la plage étudiée, l'écart reste positif et du même ordre de grandeur.
Le sens du biais d'additivité reste inconnu : mutualiser la remédiation réduit les coûts,
alors que saturer la capacité de réponse les augmente.

Un modèle de coût
à contrainte de capacité ferait du sens de l'effet une conséquence plutôt qu'un scénario. L'enjeu
se mesure : élasticités de $0{,}95$ à l'état non conforme contre $0{,}39$ à l'état conforme, soit un
facteur $2{,}42$, car l'exposant ne pèse que sur les sinistres multi-piliers. Il faudrait un registre
interne donnant le coût de chaque incident et le nombre de domaines touchés.

% SOURCES-SCRIPTS: 74

\section{La contagion entre entités, par le prestataire partagé}

La cascade court entre les cinq piliers d'une entité, non entre entités (chapitre~\ref{ch:positionnement}).
Le canal d'accumulation est pourtant ancré sur ce que les entités \emph{partagent} : la concentration
du marché de l'informatique en nuage et les prestataires critiques. Une cascade à deux étages, entre
piliers puis entre entités, donnerait un besoin de capital de place. Le registre d'informations exigé
par le règlement liste les prestataires de chaque entité : en croiser les registres donne un graphe de
fournisseurs partagés \emph{observé}, sans identifier pour autant les probabilités de transmission. Mais l'énumération des sommets
deviendrait impossible, et l'identification partielle devrait passer par une optimisation.


% =====================================================================
